\documentclass[a4paper,12pt]{report}
%\usepackage{algorithmic}
%\usepackage[linesnumbered,ruled,vlined]{algorithm2e}
\usepackage[margin=2cm]{geometry}
\usepackage[utf8]{inputenc}
\usepackage{listings} 
\usepackage{graphicx} 
\usepackage{color}
\usepackage{xcolor}
\usepackage{hyperref}
%\usepackage{mdframed}

\newcommand{\currentdata}{ 1 February 2017}
\newtheorem{example}{Example}

\begin{document}
\vspace{-5cm}
\begin{center}
Departamentul de Calculatoare \c si Tehnologia Informa\c tiei\\
Universitatea Tehnic\u a din Cluj-Napoca\\
\includegraphics[width=10cm]{fig/footer}
\end{center}
\vspace{1cm}
%\maketitle
\begin{center}
  \begin{Large}
    \textbf{Sisteme Inteligente}\\
  \end{Large}
  \textit{Activitate de laborator \^in 2016-2017}\\
\vspace{3cm}
Project title: ....\\
Tool: \\
\vspace{1.5cm}
Name:\\
Group:\\
Email:\\
\vspace{6cm}
Assoc. Prof. dr. eng. Adrian Groza\\
Adrian.Groza@cs.utcluj.ro\\
\vspace{1cm}
\includegraphics[width=10cm]{fig/footer}
\end{center}

\tableofcontents

\chapter{AI projects and tools ($W_1$)}


\colorbox{blue!20}{\fbox{\begin{minipage}{1.0\textwidth}
The teaching objectives for this week are:
\begin{enumerate}
 \item To identify existing projects for undergraduate level in the AI domain.
\item To get awareness of the effort expected for the laboratory activity by browsing past projects at TUCN and other universities.
\item To get used with the technical instrumentation used during this semester. 
%That is: i) assigning the AI tool, ii) set your latex environment for writing your final report.
\end{enumerate}
\end{minipage}}}
\vspace{0.2cm}

At the beginning of each class, please write what do you expect to learn/achieve/understand/develop during the following two hours.

\vspace{0.2cm}

\colorbox{blue!20}{\fbox{\begin{minipage}{16cm}
My personal objectives for this class are:
\begin{enumerate}
 \item 
\end{enumerate}
\end{minipage}}}
\vspace{0.2cm}
This laboratory follows a project-based learning methodology for learning artificial intelligence.

An ideal project should be one that demonstrates some creativity, 
attempts to answer an interesting research question, 
offers an interesting AI solution to a real scenario and validates the solution through graphs and experiments.
These aims cannot be achieved without reading relevant AI articles for your scenario.

\section{AI undergraduate projects}
This section aims to provide you starting ideas on student AI projects.
Browse the following resources:  
\begin{enumerate}
\item https://www.quora.com/What-are-basic-artificial-intelligence-projects-for-beginners
 \item \href{http://www.akira.ruc.dk/~keld/teaching/projektforslag/ai_projects.pdf}{Project proposals in AI} at Roskilde University, Denmark
\item \href{http://www.cs.cmu.edu/~./10701/projects.html}{Machine learning project} suggestions at School of Computer Science, Carnegie Mellon University (http://www.cs.cmu.edu/~epxing/Class/10701/project.html)  
\item Challenges in artificial intelligence domain at \href{https://www.hackerrank.com/domains}{HacckerRank}
\item Repository of AI assignments at \href{http://modelai.gettysburg.edu/}{AI-repository}
\item Collection of projects at MIT open coursware for the Knowledge-Based Applications Systems class
\href{http://ocw.mit.edu/courses/electrical-engineering-and-computer-science/6-871-knowledge-based-applications-systems-spring-2005/projects/}{MITOpenCoursware}
\end{enumerate}

The projects are intended to let you look in depth at some area of artificial intelligence 
that may only be covered briefly in class or~\cite{russell2003artificial}. 
You can see the AI lab as an opportunity for you to explore an interesting problem of your choice in the context of a real-world scenario.
Hence, we encourage you to do some original research in the AI domain that is of interest to you. 
To give you an idea, about how your fellows work as students, take a look at the examples below:

%Here is a list of project ideas:

\begin{enumerate}
 \item A machine learning approach for identifying patterns at Eurovision musical competition. 
\item Performance comparison of AI algorithms in the Wumpus world.
\item A multi-agent system for playing the board game Risk. Assess the
relative strength of each player and the strategic value of each country. 
\item Intersection situation awareness and normative reasoning.
\item Summarize an article written in wikipedia or other technical text.
\item Checking if a small text is entailed by a larger text.
\item Extracting arguments in a persuasive or legal text.
\item A knowledge-based computer purchasing adviser.
\item Build an ontology from a collection of documents using machine learning.
\item Legal assistant in case of divorces. Splitting marital property between spouse.
\item Advisor for installing a wind turbine.
\item Compliance checking of architectural plans for buildings.
\item Knowledge based system for automobile fault diagnosis.
\item Real estate analyzer system. 
\item Fake opinion detection system. 
\end{enumerate}


%Sample projects at TUCN are:

%\begin{enumerate}
% \item 
%\end{enumerate}


Projects at other universities:
\begin{enumerate}
 \item A project submitted by Greg Barish entitled "Approaches to integrating abstractions in graphplan-based planning systems" for the Artificial Intelligence Planning class at University of Southern California.
\href{http://www.isi.edu/~blythe/cs541/Projects/barish.pdf}{Final Report}
\item A project submitted by Victor L. Williamson entitled ``What to do With a Patient Who Has
Chest Pain?'' \href{http://ocw.mit.edu/courses/electrical-engineering-and-computer-science/6-871-knowledge-based-applications-systems-spring-2005/projects/william_cp_ches.pdf}{Final report}
\end{enumerate}

You may also check AI books like "Programming Collective Intelligence" by Toby Segaran~\cite{book:1527985}. 
It contains many interesting examples.

\section{AI-related competitions}
As part of your laboratory work, we encourage you to participate to various students competitions in AI:

\begin{enumerate}
\item Machine learning competitions: \href{https://www.kaggle.com/}{Kaggle}
\item Ontology development competition: \href{http://fois2014.inf.ufes.br/p/call-for-ontology-competition.html}{Ontology Competition}
\item Competition on computational models of argumentation: \href{http://argumentationcompetition.org/}{ICCMA}
\item Competitive programming \href{https://www.hackerrank.com/domains}{HacckerRank}
\item Student StarCraft AI tournament 2016 \href{http://sscaitournament.com/}{SSCAIT}
\item Power trading agent competition \href{http://www.powertac.org/node/39}{PTAC}
\item Ontology alignment evaluation initiative \href{http://oaei.ontologymatching.org/2015/}{OAEI2015}
\end{enumerate}

\paragraph{Disseminating your work.}
 Presenting your results at student conferences increases your chances to obtain a master scholarship. 
 One possible student conference is:

\begin{enumerate}
 \item Computer Science Students Conference at UBB and TUCN: 
 \href{http://cs.utcluj.ro/csd/site/index.html}{CSSC}
\end{enumerate}


\section{Running latex}
LATEX is a typesetting system suitable for producing scientific and mathematical documents. 
Three starting points are: 
\begin{itemize}
 \item \href{www.maths.tcd.ie/~dwilkins/LaTeXPrimer/}{Getting started with LaTeX (David R. Wilkins)} 
\item \href{frodo.elon.edu/tutorial/tutorial/tutorial.html}{LaTeX Tutorial} (Jeff Clark)
\item \href{http://cs-gw.utcluj.ro/~srazvan/articleSchema.tgz}{Very Brief Introduction to Latex} by Radu Slavescu
\end{itemize}


Kile is one latex editor. 
If installed, just type \texttt{kile} in the terminal. 


\section{Linux support}

Becoming familiar with the Linux requires patience. 
You must have the desire to try and figure things out on your own, rather than having everything done for you. 
Two starting references are:
\begin{itemize}
 \item \href{http://tldp.org/LDP/intro-linux/html/index.html}{Introduction to Linux A Hands on Guide, Machtelt Garrels} 
 \item \href{http://cs-gw.utcluj.ro/~srazvan/articleSchema.tgz}{Brief Synopsis of Linux} by Radu Slavescu
\end{itemize}

Table~\ref{tab:commands} lists basic commands.


\begin{table}
\caption{Quickstart linux commands.}
\label{tab:commands}
\begin{center}
\begin{tabular}{|l|l|}\hline
Command & Meaning \\ \hline
pwd & display present working directory \\ \hline
ls & displays the files in the current working directory\\ \hline
cd & change the directoris \\ \hline
chmod +x & set a file as executable \\ \hline
man & read man pages of a command\\ \hline
ssh & connects to a secure shell\\ \hline
\end{tabular}
\end{center}
\end{table}




\section{Exercises}
\begin{enumerate}
 \item Compile the \texttt{is.tex} file in order to start writing your notes. 
Recall that this documentation is also a \textit{support for you}, during the design and implementation of you ideas. 
Make an habit in writing down your ideas from the first week in a professional manner.
\item Identify 3 Web resources with ideas on student AI projects.
\item Think at one of your hobbies. Would be possible to develop something on that line?
\item Identify a media source for AI-news. Investigate interesting ideas in that news. 
Do the AI-technologies behind these ideas appear in the AIMA book?
\item Identify AI journals in Science Direct and Springer Verlag. Browse some abstracts from these journals.
\item Identify an AI competition. Consider making a team of 2-3 students to participate at that competition.
\item Imagine that you are the founder of a start-up. What kind of innovative project would be feasible for you company? 
\item Write a short list (3 to 5) of possible projects for this lab. Be ambitios. 
\item Display network information for your workstation.
\item Connect via \texttt{ssh} to another workstation. 
\end{enumerate}

\fbox{\begin{minipage}{16cm}
Solution to exercise 2       
\end{minipage}}
\vspace{0.5cm}

\fbox{\begin{minipage}{16cm}
Solution to exercise 3       
\end{minipage}}
\vspace{0.5cm}

\fbox{\begin{minipage}{16cm}
Solution to exercise 4       
\end{minipage}}
\vspace{0.5cm}

\fbox{\begin{minipage}{16cm}
Solution to exercise 5       
\end{minipage}}
\vspace{0.5cm}

\fbox{\begin{minipage}{16cm}
Solution to exercise 6       
\end{minipage}}
\vspace{0.5cm}

\fbox{\begin{minipage}{16cm}
Solution to exercise 7       
\end{minipage}}
\vspace{0.5cm}

\fbox{\begin{minipage}{16cm}
Solution to exercise 8       
\end{minipage}}
\vspace{0.5cm}

\fbox{\begin{minipage}{16cm}
Solution to exercise 9       
\end{minipage}}
\vspace{0.5cm}

\fbox{\begin{minipage}{16cm}
Solution to exercise 10       
\end{minipage}}
\vspace{0.5cm}

\chapter{Installing the tool ($W_2$)}
\colorbox{blue!20}{\fbox{\begin{minipage}{16cm}The teaching objectives for this week are:
\begin{enumerate}
 \item To install the tool on a Linux system
\item To get aware of tool plugins, extensions, and running parameters.
\end{enumerate}
\end{minipage}}}

\vspace{0.5cm}

\colorbox{blue!20}{\fbox{\begin{minipage}{1.0\textwidth}
My personal objectives for this class are:
\begin{enumerate}
 \item 
 \item 
\end{enumerate}
\end{minipage}}}
\vspace{0.2cm}

The AI lab follows the problem-based learning model~\cite{hmelo2004problem}. 
In this model, you have to:
\begin{enumerate}
 \item identify a real world situation that has no clear (or right, or efficient) 
solution;
\item  develop a viable solution to the identified problem;
\item  get new information through self-directing learning;
\item  take the decisions needed during the lifecycle of your project 

\end{enumerate}


%(staff act as faccilitators by helping you during the cognitive process of reflection)
The first sub-problem that you are facing is to install 
the tool on a Linux distribution.
%Linux advantages

First, you need to find the version of your operating system.
To find the kernel version, type in a terminal \texttt{uname -a}.
To find the Linux distribution type \texttt{lsb\_release -a}.

Then you need to download the distribution of the tool adequate for your system.

Extract the archive in your own directory. 
Avoid using spaces for the directory name. 
Identify a README file or Installation instructions for the tool.

When browsing the requirements, check first whether they have been not 
already installed.
For instance, Java might be already in the system. 
You can check this with \texttt{java -version}.

List here the steps done for installing the tool. 
A typical installation could look like:

\begin{verbatim}
./configure
make
make install
export PATH=PATH:......
\end{verbatim}

List here the exact commands for running the tool. 
These lines will save you precious time in the next week.

Finally, you need to check the installation.


In case something goes wrong:

\begin{enumerate}
 \item Try to install a newer/older version of the tool
\item Try to obtain the source code and compile it.
\item Install the proper version for the software requirements (i.e., correct LISP or Prolog),
\item Take a look into the installation/running scripts or in the \texttt{makefile} file.
\item Try to identify a version of the tool written in other programming language.
\end{enumerate}

It is your interest not to get stuck with this installation. 
It is therefore important to identify a solution such that you 
can start shaping your project from the next laboratory. 


\section{Exercises}
\begin{enumerate}
 \item List the steps done for installing the tool. 
 \item List the exact commands needed to run the tool. 
 \item What other AI tools can use the output of your tool? 
 Can you indentify such tools on the Web?
 Try to assess the difficulty level and risks of integrating such tools for your project.
\end{enumerate}

\fbox{\begin{minipage}{16cm}
Solution to exercise 1       
\end{minipage}}
\vspace{0.5cm}

\fbox{\begin{minipage}{16cm}
Solution to exercise 2       
\end{minipage}}
\vspace{0.5cm}

\fbox{\begin{minipage}{16cm}
Solution to exercise 3       
\end{minipage}}
\vspace{0.5cm}




\chapter{Running and understanding examples ($W_3$)}

\colorbox{blue!20}{\fbox{\begin{minipage}{1.0\textwidth} 
                                          
The teaching objectives for this week are:
\begin{enumerate}
 \item To run and understand the toy examples provided by the tool
\item To identify what realistic problems are adequate for your tool 
\end{enumerate}
\end{minipage}}}

\vspace{0.5cm}

\colorbox{blue!20}{\fbox{\begin{minipage}{1.0\textwidth} 
My personal objectives for this class are:
\begin{enumerate}
 \item 
 \item 
\end{enumerate}
\end{minipage}}}


Describe the examples that you run. 
What type of problems can be supported by the tool?
 
 \section{Exercises}
\begin{enumerate}
 \item Detail one example that you run. Which are the input and the output? 
 Describe the structure of the code. 
 \item Describe the real world problems that can be solved by your tool/algorithm.
\end{enumerate}

\fbox{\begin{minipage}{16cm}
Solution to exercise 1       
\end{minipage}}
\vspace{0.5cm}

\fbox{\begin{minipage}{16cm}
Solution to exercise 2       
\end{minipage}}
\vspace{0.5cm}

 
\chapter{Understanding conceptual instrumentation ($W_4$)}
\label{ch:tool}


The teaching objectives for this week are:
\begin{enumerate}
 \item To understand the algorithm(s) on which your tool relies.
\item To get used with writing algorithms in Latex
\end{enumerate}

Latex provides various environments for writing algorihms, including: 
\texttt{algorithm, algorithmic, algorithmicx, algpseudocode, algorithm2e}.

The following example algorithm has been taken from~\cite{7034697}.


\section{Exercises}
\begin{enumerate}
 \item Big O complexity
 \item Which are the latex options to write algorithms? 
 Describe in one paragraph the main features of one such package for algorithms.
\end{enumerate}

\fbox{\begin{minipage}{16cm}
Solution to exercise 1       
\end{minipage}}
\vspace{0.5cm}

\fbox{\begin{minipage}{16cm}
Solution to exercise 2       
\end{minipage}}
\vspace{0.5cm}



\chapter{Project description ($W_5$)}

The teaching objectives for this week are:
\begin{enumerate}
 \item To have a clear description of what you intend to develop.
\item To point to specific resources (datasets, knowledge bases, external tools) 
that support the development of your idea and which minimise the risk of failure.
\item To identify related work (articles) that are relevant or similar to your approach.
\end{enumerate}

\vspace{0.5cm}

My personal objectives for this class are:
\begin{enumerate}
 \item 
 \item 
\end{enumerate}


To encourage the development of AI skills, 
you  were required to come up with a significant semester project 
You have to  apply ideas from the course to a problem of your own.

Which domain to choose is a decision that only you can make. 
The more aware of the tool capabilities, the more adequate the decision.
Realistic and original scenarios are encouraged. 
Well known toy problems (salesmen, map colouring, logistic planning, wumpus, sudoku, queens, missionaries and canibals, etc.) do not worth much for your grade. 
Your scenario should be realistic and should be business oriented. 

Select clearly defined problems and not generic ones 
(i.e I will do something in the medical domain). 
Note that the focus is both on programming and on modelling the reality into a formal representation.


Before specifying your project you must understand as much as possible about the application domain (medicine, bank, human resource management, etc). 
You must also understand functionality required by the stakeholders of your system. 
Let the problem drive the modelling - the more you understand the domain, the more technical solutions need to be solved by you. 

Consider answering to the following questions:
\begin{enumerate}
 \item What will your system do?
\item Which is the scope of coverage your system aims for?
\item What will be the input of your program?
\item What will be the output of your program?
\item What will be the knowledge of your system?
\item Which would be the narrative description of running scenario(s)?
\item Which are the stakeholders of your system?
\item Which are the asumptions?
\end{enumerate}

\begin{example}[What will system do]
 
\end{example}

\begin{example}[Scope of the program]
 
\end{example}

In this problem-based model you learn what you 
need to know in order to solve a problem. 
However, note that you have total flexibility in stating your project objectives.
For instance, if you consider that studying more 
than one computational technology (or AI algorithms) brings more benefits, 
you are encourage to do it. 
You just have to frame your task under one scenario umbrella. 
One example: you can investigate the problem of {\it fake review detection} with 
various machine learning algorithms: decision trees, nayve bayes, neural networks, ensemble learning. 
This road helps you to study and compare the above algorihtms on your own problem.
A second example: you can help a {\it robot to escape from a maze} 
with various search algorithms (A*, greedy, deep first, uniform cost) 
or computation technologies (constraint satisfaction problems, planning, 
searching with observations).
Hence, both the problem-based model or a more algorithmic approach to AI 
(if formulated under the umbrella of a single scenario) 
are accepted for this laboratory.



\section{Narrative description}
This should be a simple textual description of your scenario.
You should explain your project objectives.
Put them all together in half a page.

\section{Facts}
You start the analyze of the problem by identifying 
the relevant facts from the scenario.
This fact-identification step helps you to represent the problem.


\section{Specifications}
List of specifications. 
Use your knowledge from "System Engineering" on how to write specifications and requirments (see Exercise 1).

\section{Top level design of the scenario}
Technical description of your scenario. 
In some cases, you may provide a figure with the architecture of the system. 


\section{Knowledge acquisition}
You should be aware that computing interacts with many different domains. 
Solutions to many AI problems require both computing skills and domain knowledge.

First, you should ask yourself if you have the necessary background and resources to do
a project in the chosen area. 

\paragraph{How do represent knowledge?} 
Your system relies on a knowledge base.
You have to describe how do you represent this knowledge.
You might choose between different logics: 
propositional logic, first order logic, modal logics, description logics, epistemic logics, 
temporal logics, and so on. 

\paragraph{Where are you getting the required knowledge/data}
Point towards the knowledge bases that you plan to exploit. 
The existence of these sources are required to prove that your approach is realistic.

Examples of knowledge sources include:
\begin{itemize}
\item Data sets: i.e., https://archive.ics.uci.edu/ml/datasets.html 
\item Statistics: i.e., http://ec.europa.eu/eurostat 
\item Ontology repositories in OWL or RDF format.
\end{itemize}


If you will be using books, give their reference. 
If you hope to exploit people for elicitation, give their names.
If you aim to use data sources or knowledge repositories list them and be sure that you have 
access to the needed knowledge. 
Indicate what you have accomplished so far in knowledge acquisition. 




\section{Related work}
You have to identify articles or conference papers relevant to your scenario.
Searching for adequate references can be both rewarding and frustrating.

Browse online libraries like:
\begin{itemize}
 \item Science Direct
 \item IEEE Computer Society Digital Library: IEEEXplore
 \item SpringerLink:  http://www.springerlink.com/computer-science/
 \item ACM Digital library: http://portal.acm.org/dl.cfm
%\item The Collection of Computer Science Bibliographies: http://liinwww.ira.uka.de/bibliography/
%\item Computer Science Technical Reports: http://www.nzdl.org/
\end{itemize}
A valuable resource is Google Scholar. 
Some references may be frely available on ResearchGate. 

Looking at the examples remains the best ways of learning how to present a literature analysis. 
Eacy article does include such section.


Obtain the \texttt{.bib} file of each article that you will rely on. 
Cite and very briefly describe the main idea of each paper that you have read.
Save the most relevant related papers in a local directory.
Use your own words. 
Don't use automatic translation tools (i.e., Google translate) just to expand your documentation.
You should already be aware that this is a form of cheating. 
Everything in your project that does not come with a citation is assumed to be your own work.


The following is an example of such \texttt{bib} structure:

\begin{verbatim}
 @article{bench-capon:argumentation-in-ai,
 author = {Bench-Capon, Trevor J. M.  and Dunne, Paul E. },
 title = {{A}rgumentation in {A}rtificial {I}ntelligence},
 journal = {Artificial Intelligence},
 volume = {171},
 number = {10-15},
 year = {2007},
 issn = {0004-3702},
 pages = {619--641},
 doi = {http://dx.doi.org/10.1016/j.artint.2007.05.001},
 publisher = {Elsevier Science Publishers Ltd.},
 address = {Essex, UK}
}
\end{verbatim}

Don't forget to include the above structure in your \texttt{.bib} file. 
Then, generate the \texttt{.bbl} file in order to correctly appear in the \textit{Bibliography} section.

\section{Exercises}
\begin{enumerate}
\item Sum up what is the aim of your project in a Twitter-sized phrase (140 characters)
\item Recall or identify an engineering methodology to write specifications. 
 Cite this methodology and employ it for specifying your project.
 \item Which are the differences between requirments and specifications? 
 \item Identify similar scenarios proposed by your collegues. 
 Think at some form of collaboration with one of your collegues having similar interests.
\end{enumerate}

\fbox{\begin{minipage}{16cm}
Solution to exercise 1       
\end{minipage}}
\vspace{0.5cm}

\fbox{\begin{minipage}{16cm}
Solution to exercise 2       
\end{minipage}}
\vspace{0.5cm}

\fbox{\begin{minipage}{16cm}
Solution to exercise 3       
\end{minipage}}
\vspace{0.5cm}

\fbox{\begin{minipage}{16cm}
Solution to exercise 4       
\end{minipage}}
\vspace{0.5cm}



\chapter{Preliminary results ($W_7$)}
This section corresponds to the midway report in week 7.
The teaching objectives for this week are:
\begin{enumerate}
 \item To prove that you have managed to write few lines of code of your own.
\item To prove that the knowledge or data required are already obtained.
\end{enumerate}

These objectives decreases the risk to fail. 
You should be aware that failing to meet the above objectives in week 
7 indicates high risks in obtaining relevant results at the end of the semester.
Take urgent measures to overcome these difficulties.


\section{Exercises}
\begin{enumerate}
\item Write the preliminary results explaining any realizations or insights found during the research of the subject.
\item Discuss new information and questions found during the domain investigation or during coding.
\end{enumerate}








\chapter{Implementation details ($W_9$)}

The teaching objectives for this week are:
\begin{enumerate}
 \item Illustrate each aspect of the reality that you have 
 modelled in your solution.
\item To explain the relevant code from your scenario.
\end{enumerate}

\vspace{0.5cm}

My personal objectives for this class are:
\begin{enumerate}
 \item 
 \item 
\end{enumerate}

Projects in artificial intelligence consist of developing new solutions. 


\section{Relevant code}

Provide the relevant code (see an example in Fig.~\ref{fig:code}).
You can use "verbatim" package or "listing" package. 
Complement the code with its corresponding textual description.

\begin{figure}
\begin{verbatim}
(full-reset)
(instance a Argument)
(related a b attacks)
(concept-instances Argument) 
\end{verbatim}
\caption{Modelling arguments in Racer.}
\label{fig:code} 
\end{figure}

The eager student may use concepts from \textit{literate programming}.

\section{Common bad practice in AI undergraduate projects}

\paragraph{The over-estimated AI programmer.} $ $

{\it Bad practice}: Excepting few genial students, you tend to overestimate your AI-programming abilities. 
That is, you start to write a large amount of code. (Here large might be 20 lines).
When testing it, nothing run. 
You start to debug a line or to remove it.
Your program will not run this time too. 
You remove or comment another line. 
And so on, until you have a single line of code.
If you are lucky, that could run. 
But you lose a lot of time in this enterprise.

{\it Solution}: In the early stage of writing code, write a line of code and test it. 
If it works, write another line and test it. 
And so on. 
That is, you are exploiting the interactive environments provided by 
AI tools or languages like LISP and PROLOG.
You should hold your horses and 
have the most possible skeptical attitude towards your code.  
As you get experience, your will be noticing that writing 
AI-declarative code is more effective than procedural one.  


\paragraph{The eyewash bug.} $ $
{\it Bad practice}: You spend most of your programming time to develop a GUI for your AI-system. 
Don't bother. I am sympathetic with Sania Twain's view on GUIs: "You don't impress me much".
Such things are indeed important in computer science, but not relevant in this AI class.


%\paragraph{The stucker bug} 

\paragraph{The not-organised student.}
You are not organised, if something like this will happen to you:
\begin{itemize}
 \item You do not find your project and yield "Someone removed my project!". 
Most of the time your are logged with a different user as usual. 
Check this with \texttt{who am i}. 
This is not a rhetorical question, but a Linux command.
\item You are working in a different directory. 
Type \texttt{pwd} and \texttt{ls} to check that your executables are indeed in the current working directory.
If you have been lazy to set your PATH variable, you might just forgot to type \texttt{./} 
for executing the command in the current directory.
\end{itemize}

\paragraph{The omniscient student.}
You are in this cathegory if you fail to add references. 
Reading relevant references is mandatory to deliver a decent project. 

\section{Exercises}
\begin{enumerate}
 \item What latex packages can be used to format code? 
 \item 
\end{enumerate}

\fbox{\begin{minipage}{16cm}
Solution to exercise 1       
\end{minipage}}
\vspace{0.5cm}

\fbox{\begin{minipage}{16cm}
Solution to exercise 2       
\end{minipage}}
\vspace{0.5cm}




\chapter{Tool expressivity ($W_{10}$)}


The teaching objectives for this week are:
\begin{enumerate}
 \item Describe each technical instrumentation provided 
by the tool that was enacted in your implementation.
\end{enumerate}

\vspace{0.5cm}

My personal objectives for this class are:
\begin{enumerate}
 \item 
 \item 
\end{enumerate}



%For instance:
%{\it "Racer offers also the capability to use {\it rules}. 
%he code in Fig.~\ref{fig:rules} illustrates my rule used to create instances of type student."}


\chapter{Graphs and experiments ($W_{11}$)}

The objectives for this week are:
\begin{enumerate}
 \item To describe and interpret each experiment that you have performed
\end{enumerate}

\vspace{0.5cm}

My personal objectives for this class are:
\begin{enumerate}
 \item 
 \item 
\end{enumerate}


An experiment investigates how some variables are related. 
Usually, experiments verify a previosly formulated hypothesis.
Such hypothesis may investigate how your software degrades its performance with larger inputs.
You will need to run simulations to see how your implementation is affected by different inputs.


Note that running experiments mean more than testing your solution. 
It helps to describe and prove how did you test your implementation.
Moreover, during this lab, you will often need to: 
1) generate random data for your algorithms,
2) measure their performance (number of operations, execution time),
3) draw charts,
4) interpret the obtained results.

The eager student might want to take a look at literature on how to design computer experiments, such as~\cite{fang2005design}. Section 5.6 from~\cite{fang2005design} might be of particular interest for some of you. 
If your experiments include a stochastic parameter, you need to include a test for statistical significance. 
This is important to prove that your ouputs are not a random effect.



\begin{figure}
%\includegraphics[width=8cm]{} 
Graphs always impress teachers...
\caption{Increasing the accuracy with the number of samples.}
\label{fig:accuracy}
\end{figure}

You should develop a test suite that can be used to show your code works correctly under a various conditions/problems/scenarious.


\section{Evaluation metrics}



\chapter{Related work and documentation ($W_{12}$)}

The teaching objectives for this week are:
\begin{enumerate}
 \item To compare your results to related work.
\item To discuss the advantages and limitations of your solution.
\item To deliver a professional documentation of your work.
\end{enumerate}

\vspace{0.5cm}

My personal objectives for this class are:
\begin{enumerate}
 \item 
 \item 
\end{enumerate}


This chapter convinces me that you know how your work fits into the larger domain area.


%For instance: \cite{chesnevar:Survey2000}

\section{Related approaches}
You need to support your opinions with trustworthy evidence and references.
You need also to decide how your problem fits into a wider context.
The quality of your reference is a strong indicator that you 
managed to scrutinise different perspectives on the topic.
Proving understanding of the references is the foundation of a good grade. 
By start coding without reading relevant references, you will most probable write something irelvant for the application domain.  

Identify and describe other solutions for solving the same (or similar) scenario like yours.

\section{Advantages and limitations of your solution}
This part of the conclusions chapter should be an evaluation of your work. 


\section{Possible extensions of the current work}






\chapter{Project demo and documentation ($W_{13}$)}

The teaching objectives for this week are:
\begin{enumerate}
\item Deliver the technical documentation of your project
\item Demonstrate your running scenario to the instructor
\end{enumerate}

\vspace{0.5cm}

My personal objectives for this class are:
\begin{enumerate}
 \item 
 \item 
\end{enumerate}


Demonstrate in 4-5 minutes your running scenario to the instructor. 
The demo should take place on a Linux distribution


From your final report, remove the text/examples/algorithms/rules/bibliographic references/etc - keep only your notes. 
If the documentation does not meet minimum standard for lisability and scientific discourse, 
it will be classified by the furious teaching assistant as unacceptable and therefore rejected.




\section{Exercises}
\begin{enumerate}
 \item How you would advocate your project to a possible client?
 \item Write five highlights of your results (maximum 85 characters including spaces).
 \item 
\end{enumerate}

\fbox{\begin{minipage}{16cm}
Solution to exercise 1       
\end{minipage}}
\vspace{0.5cm}

\fbox{\begin{minipage}{16cm}
Solution to exercise 2
\end{minipage}}
\vspace{0.5cm}


\chapter{Results dissemination and feedback ($W_{14}$)}

The teaching objectives for this week are:
\begin{enumerate}
\item To practice public presentation.
\item To get used with the beamer template for making scientific presentations
 \item To get feedback from your collegues. 
\end{enumerate}

\vspace{0.5cm}

My personal objectives for this class are:
\begin{enumerate}
 \item 
 \item 
\end{enumerate}

Learning is enhaced by constructive feedback on the strong/weak points of your performance during AI laboratory. 
This feedback focuses on the scientific relevance of your results and it aims to complement the feedback encapsulated in the grade.
Do not take criticism personally. 
It is the project that is being critisised, and not your competence or intelligence. 


\subsection{Public presentation}
10 slides in beamer format for a timeslot of 5 minutes presentation plus 5 minutes questions.
Questions may be posed by your collegues or the teacher. 

One of the main difficulties when designed slides is how to find the right amount of 
technical details to be included. 
No technical details rise the question of "bla bla story telling". 
To much technical details may bore the audience and also you may fail to fit within the time assigned. 

Two introductory tutorials on beamer are~\cite{mertz2005beamer} and~\cite{batts2007beamer}.
The beamer manual is: 

\href{http://cs-gw.utcluj.ro/~srazvan/articleSchema.tgz}{Presentation template}

During presentation, it is a mistake to focus on the tool that you have been used. 
During 5 minutes, you have to focus only on your results and to market your work.
Don't forget to include technical details and graphs.

A method for disseminating yor research consists of writing 3-5 highlights.
Highlights consist of a set of bullet points 
that convey the core findings of your work.
For examples, see http://www.elsevier.com/highlights.

You are now playing the role of your project advocator. 
Always keep in mind that you have to market your project only, and not the tool that you rhave relied on.



\subsection{Self-assessment}

\begin{table}
\begin{tabular}{p{8cm}l}
Aspect & Self-assessment\\ \hline
How did you manage to master the tool?& \\
How realistic was your scenario? & \\
Relevance of the running experiments &\\
Knowledge and skills achieved & \\ 
Capacity to market your effort and results through documentation and presentation& \\
\hline
\end{tabular}
\caption{Self-assessment. Assess each aspect, with: enthusiastic, satisfacatory, unsatisfacatory, bad} 
\end{table}

‘What did you do well? Give examples’

‘Where do you think the assignment is weak?’


\subsection{Formative feedback}
The last week is an opportunity for interested students to obtain 
a formative feedback. 
This is not an opportunity to negatiate your grade. 
In the previous week you had your chanches to advocate your work. 



\subsection{Problem-based learning}

One scope was to engage you in a kind of self-directed learning.
The rationale is that you are more heterogeneous and 
you have different learning experiences and maturity levels.
The focus was not on mastering AI algorithms but to apply them in practice. 
By practice I mean realistic scenarious.
Ideally, you should have developed more awareness of the social, 
environmental, economic aspects of a real problem. 


\appendix

\chapter{Your original code}
\label{app:code}
This section should contain only code developed by you, without any line re-used from other sources. 
This section helps me to correctly evaluate your amount of work and results obtained. 
Including in this section any line of code taken from someone else leads to failure of IS class this year.
Failing or forgetting to add your code in this appendix leads to grade 1.
Don't remove the above lines.




\chapter{Quick technical guide for running your project}

Requirments

Step by step technical manual

%\input{mycode}

\chapter{Check list}

\begin{enumerate}
 \item Your original code is included in the Appendix~\label{app:code}.
 \item Your original code and figures are readable.
 \item All the references are added in the Bibliography section.
  \item All your figures are referred in text (with command ref{}), described in the text, and they have relevant caption.
 \item The final documentation describes only your project. Don't forget to remove all tutorial lines in the template (like these one).
 \item The main algorithm of your tool is formalised in latex in chapter~\ref{ch:tool}.
  
\end{enumerate}


\bibliographystyle{plain}
\bibliography{is}


\vspace{2cm}
\begin{center}
Intelligent Systems Group\\
\includegraphics[width=10cm]{fig/footer}
\end{center}



\end{document}
